\section{Related Work}
\label{sec:related}

\subsection{Safety Evaluation Methodology}

Benchmarks address bias~\cite{parrish2022bbq, gallegos2024bias}, truthfulness~\cite{lin2022truthfulqa}, over-refusal~\cite{rottger2024xstest, cui2025orbench}, and sycophancy~\cite{perez2023discovering, sharma2024sycophancy}, each operationalising its target construct differently enough that cross-benchmark comparison is not on offer, and rigour has not kept pace with proliferation.  Safety benchmark scores correlate strongly with general capabilities~\cite{ren2024safetywashing}, casting doubt on whether they measure anything safety-specific; systemic issues in safety-benchmark design have been catalogued~\cite{eriksson2025trustbenchmarks}, and 83\% of agentic-AI evaluations focus on technical rather than human-centred dimensions~\cite{jafari2025imbalance} (a ratio that flips the priority order most safety researchers would endorse).  Mou et al.~\cite{mou2024sgbench} document limited generalisation of safety alignment across task and prompt types, with most LLMs performing worse on discriminative than generative tasks; that work does not examine the MC-to-OE shift specifically, nor its interaction with deployment architecture.

Existing critiques focus on benchmark design and scoring methodology; whether the response format itself moves measured safety has gone largely unaddressed.

\subsection{Format Effects in LLM Evaluation}

Capability evaluation has a substantial literature on format sensitivity --- position bias, option-order effects, MC-to-OE accuracy drops --- but the same lens has not been turned on safety benchmarks.

\paragraph{Selection bias and option-order sensitivity.}
Zheng et al.~\cite{zheng2024selectionbias} document position bias in MC selection: LLMs systematically prefer certain letter positions over others, independent of content. Reordering options produces 13--75\% accuracy swings across benchmarks, even under few-shot conditions~\cite{pezeshkpour2024ordersensitivity}. The joint contribution of option ordering and token identity has been quantified across model families~\cite{wei2024unveilingbias}. MC evaluation is therefore not construct-neutral: it introduces systematic biases independent of what the benchmark is trying to measure.

\paragraph{The MC-to-open-ended gap.}
The gap between MC and open-ended performance is large and consistent. Myrzakhan et al.~\cite{myrzakhan2024openllm} document an average $\sim$25~pp drop from MC to open-ended across every model tested. Wang et al.~\cite{wang2024mcrobust} find that instruction-tuned models evaluated via text generation are more perturbation-resistant under MC than under first-token-probability scoring; format sensitivity depends on both the evaluation procedure and the response format. Chandak et al.~\cite{chandak2025answermatching} show that many prominent-benchmark MC questions leak enough information through the answer options that models can identify the correct choice without reading the stem at all; what such benchmarks measure is option recognition, not question answering. Wang et al.~\cite{wang2024leastincorrect} establish that LLMs often answer MC questions by elimination (choosing the least-wrong option) rather than by recognising the right one; substituting ``None of the Above'' for the correct answer reduces accuracy by up to 70.9\%. These studies converge on the same conclusion: MC format does not add random noise but induces a fundamentally different response process (recognition from a closed list) than open-ended generation, and therefore likely measures different constructs even with identical question text.

\paragraph{Standardisation efforts.}
Gu et al.~\cite{gu2024olmes} propose OLMES, a standards effort that codifies the methodological choices (prompt formatting, contextual information, probability-distribution normalisation, task definition) driving large performance differences on identical questions. Wang et al.~\cite{wang2024mmlupro} introduce MMLU-Pro by expanding choice sets from four to ten, an implicit acknowledgement that the four-option default is too easy for current frontier models. Both initiatives acknowledge how thoroughly format choices contaminate capability assessments; neither extends to safety benchmarks, where the stakes of a format artifact are categorically different (an MMLU artifact misranks models on a leaderboard; a BBQ artifact mischaracterises whether a model exhibits social bias).

\paragraph{Gap between capabilities and safety: what we fill.}
The studies above concern capability assessment. On the safety side, the analogous questions --- whether MC-to-OE format effects move measured safety, and whether deployment architecture interacts with that movement --- have not been asked. Format sensitivities documented for capability evaluation generalise to safety while imposing categorically different penalties: a capability gap misplaces models on a leaderboard, whereas a comparable safety gap can reverse the qualitative conclusion of an assessment (whether a model demonstrates bias, for instance; see Section~\ref{sec:format-dependence}).

The map-reduce scaffolding strips MC answer options during task decomposition, which inadvertently converts a multiple-choice item to an open-ended one. A scaffold that preserves option choices recovers 40--89\% of losses due to format conversion. The driver of the map-reduce result is measurement condition change, not reasoning disruption or deterioration in alignment; this shifts the central question of agentic safety evaluation from ``do scaffolds decrease safety?'' to ``do scaffolds change what safety benchmarks measure?''

\subsection{Agentic AI Safety}

A parallel literature examines safety in agentic deployments, establishing that deployment configuration causally affects safety outcomes but leaving two critical gaps: no existing work isolates scaffold architecture from other deployment variables while controlling for all confounds, and none considers whether the evaluation format of the underlying safety benchmark interacts with the scaffold architecture.

\paragraph{Scaffold design as a safety variable.}
Rosser and Foerster~\cite{rosser2025agentbreeder} achieve 79.4\% safety uplift through evolutionary search over multi-agent scaffolds; scaffold design is a first-class safety variable on their evidence. Yin et al.~\cite{yin2024safeagentbench} find that agent architecture affects safety more than model choice across 750 embodied tasks. MacDiarmid et al.~\cite{macdiarmid2025emergent} show that RLHF safety training calibrated on chat evaluations fails to prevent emergent misalignment in agentic settings, consistent with our central thesis: safety properties measured under one deployment configuration may not transfer to another. Vijayvargiya et al.~\cite{vijayvargiya2025openagentsafety} report unsafe behaviour in 51--73\% of safety-vulnerable scenarios across realistic agent deployments with browser, code-execution, and file-system access.

\paragraph{Reasoning chains.}
Jiang et al.~\cite{jiang2025safechain} find that longer reasoning chains reduce safety across 13 models, positioning chain-of-thought as a first-order safety variable. Huang et al.~\cite{huang2025safetytax} document a ``safety tax'' of 7--31\% reasoning-accuracy loss from alignment-targeted training, evidencing a capability-safety trade-off under scaffolding. System prompts shape agent behaviour as much as the underlying model~\cite{breunig2026systemprompts}, and compound AI systems produce emergent properties absent from individual components~\cite{du2024improving, chen2024scalinglaws, zaharia2024compound}.

\paragraph{Agentic safety benchmarks.}
The agent-safety benchmark literature is active. Zhang et al.~\cite{zhang2024agentsafetybench} evaluate 16 models on Agent-SafetyBench (2{,}000 test cases); no model exceeds 60\% safety. Andriushchenko et al.~\cite{andriushchenko2025agentharm} report that leading LLMs are unexpectedly compliant with explicitly malicious tasks under agent-style framing. Zhang et al.~\cite{zhang2025asb} formalise offensive and defensive mechanisms for ReAct architectures. Ma et al.~\cite{ma2025safetyatscale} and Wang et al.~\cite{wang2025safetyreasoningsurvey} provide broad surveys of the area.

\paragraph{The gaps we fill.}
No prior research compares a single model across multiple controlled scaffold conditions with all other variables held constant; nor has anyone examined how evaluation format interacts with scaffold architecture.  Our results show that the apparent map-reduce safety loss is almost entirely an artefact of format conversion rather than reasoning disruption, which moves the central question from ``do scaffolds reduce safety?'' to ``do scaffolds change what safety benchmarks measure?''  The interaction is invisible to any researcher who does not independently manipulate both scaffold architecture and evaluation format.

\subsection{Evaluation Methodology from Adjacent Fields}

We adapt methodology from fields that have confronted analogous measurement crises (pre-registration and specification curves from the replication-crisis literature; equivalence testing and blinding from clinical research) and show that these tools transfer directly to AI safety evaluation.

Pre-registration~\cite{nosek2018preregistration}---the cornerstone defence against post-hoc rationalisation---is almost absent from AI safety evaluation. Hofman et al.~\cite{hofman2023preregistration} extend the principles to predictive modelling in machine learning, demonstrating feasibility in computational settings. Ioannidis~\cite{ioannidis2005false} shows that excess analytic flexibility inflates the false-discovery rate to the point where most published findings are wrong; scaffold safety researchers have the same flexibility (scoring method, model selection, prompt design, statistical specification) and no convention restraining its use.

Specification curve analysis~\cite{simonsohn2020specification} addresses this by enumerating defensible analytic specifications and testing whether findings are robust across them. Simson et al.~\cite{simson2024onemodemanyscores} apply the technique to ML fairness, showing that model design decisions substantially affect measured fairness; the analogy to our finding (that scoring methodology can manufacture or reverse safety findings) is direct. We enumerate 29 pre-registered researcher degrees of freedom (Appendix~\ref{app:spec-curve}); the primary curve varies three of these (benchmark subset, model subset, scoring method) across 18 specifications; a broader exploratory curve varies nine across 384; the remaining seventeen are documented for full pre-registration transparency.

We have applied the same reporting principles as existing studies~\cite{schulz2010consort}: single-blind assessment (the scaffold artifacts were removed prior to scoring), transparent deviation reporting, and an intention-to-treat approach. All comparisons are evaluated against pre-registered TOST equivalence bounds of $\pm 2$~pp~\cite{schuirmann1987tost, lakens2017equivalence}, separating ``no evidence of harm'' from ``evidence of no harm''; statistically significant but TOST-equivalent effects (ReAct, RD~$=\,-0.7$~pp) are reported as both, and current scaffold safety evaluations almost never make either distinction. The ICH E9 guidelines~\cite{iche9} inform our approach to multiplicity correction and sensitivity analysis.

The format-dependence finding itself carries methodological weight. The 40--89\% recovery through an option-preserving variant, replicated across 384 analytic specifications, establishes format conversion (not reasoning disruption) as the dominant mechanism, with robustness across scoring method, model subset, and statistical specification. Psychology learned through the replication crisis~\cite{opensciencecollab2015replication} that under-determination is a structural problem rather than a per-paper one. AI safety evaluation sits in the same epistemic position, and the methodology assembled here is designed against it.

\subsection{LLM-as-Judge}

LLM-as-judge is now the default method for assessing open-ended model output. Zheng et al.~\cite{zheng2023judging} introduce MT-Bench and Chatbot Arena and report that frontier LLMs reach above-80\% agreement with human evaluators on open-ended outputs, while exhibiting positional, verbosity, and self-enhancement biases. Gu et al.~\cite{gu2024surveyllmjudge} survey the design space for reliable LLM-as-judge systems. Ye et al.~\cite{ye2025biasllmjudge} decompose these biases at three levels of analysis: judge, candidate, and task. Shankar et al.~\cite{shankar2024validates} document ``criteria drift'' in LLM evaluation pipelines: the criteria themselves shift over the course of a single assessment run.

We have used our tiered scoring strategy because we know these risks. Three of the four metrics score via an automatically deterministic process (multiple choice answers extracted from ground-truth keys), thus they are impervious to judge bias. The fourth metric (XSTest/OR-Bench) was scored using LLM-as-judge methodology (Gemini~3~Flash primary, Claude Opus~4.6 validation on a 10\% subsample) as the pre-registered primary method. None of the models were allowed to evaluate their own outputs (Opus is a cross-evaluator for the other models' outputs; it does not serve in this capacity with respect to its own output). The specification curve allows readers to determine whether results would be consistent across alternative scoring specifications: they are for map-reduce (100\% significant across all scoring variants) but inconsistent for multi-agent (43.5\%), a difference unobservable by those who utilize only one scoring procedure.

The format-dependence finding introduces an additional dimension: differences between MC and open-ended safety scores could reflect scoring methodology rather than model behaviour. Our 18 falsification tests, inter-judge agreement analyses, and explicit separation of format effects from scoring-method effects in the specification curve address this confound (Section~\ref{sec:format-dependence}).

Weidinger et al.~\cite{weidinger2024holistic} argue for holistic safety evaluation combining automated benchmarks with human assessment. The format-dependence finding adds urgency: if MC and open-ended formats measure different constructs, holistic evaluation must include both formats to avoid systematic blind spots.
